\section{The Schema Nobody Can Retire}
\label{sec:ch16-the-schema-nobody-can-retire}

\index{schema growth!what a missing usage record does to a schema|(}%
Section~\ref{sec:ch16-the-gate-and-what-it-needs} left the second comparison
unanswerable without a control plane: is anybody still asking for this field.
A team that cannot answer that does not stop working. They make a decision
instead, the same decision every time, and after two years the schema has a
shape.

\index{deprecation!the field this book still carries}%
This book has already made it once. \code{sessionById} was deprecated in
chapter~\ref{ch:schema-design}, with a reason string naming its replacement,
and twelve chapters later it is still in the schema and still answering.
Nothing about that is negligence. The
question that would retire it is whether anybody still sends it, and this
graph has no way to find out, so the safe answer and the lazy answer are the
same answer and I took it.

\index{schema growth!removal needs a fact nobody has}%
Add a field and no client can break, because no client selects a field it has
never heard of. Remove one and you need a fact you do not have. Two moves,
one of them free and one of them requiring evidence nobody collected, applied
by every team for as long as the graph lives.

\subsection{One Query Type, Four Teams}

\index{Query type@\code{Query} type!seven fields, four owners}%
Federation sharpens this in a way that is easy to miss, because the shared
thing does not look shared. Here is every root field in this book's
supergraph, and the subgraph that declares it:

\begin{center}
\footnotesize
\begin{tabular}{l l}
\toprule
Root field & Declared by \tabularnewline
\midrule
\code{node} & \code{sessions} and \code{speakers} \tabularnewline
\code{nodes} & \code{sessions} and \code{speakers} \tabularnewline
\code{sessions} & \code{sessions} \tabularnewline
\code{sessionById} & \code{sessions} \tabularnewline
\code{speakers} & \code{speakers} \tabularnewline
\code{ratingCount} & \code{ratings} \tabularnewline
\code{searchSessions} & \code{search} \tabularnewline
\addlinespace
\code{rescheduleSession} & \code{sessions} \tabularnewline
\bottomrule
\end{tabular}
\end{center}

\index{Query type@\code{Query} type!nobody owns it}%
Seven fields on \code{Query} and one on \code{Mutation}, from four services,
and no service owns the type they sit on. Each of those teams can add a field
to it without asking anybody, which is the thing
chapter~\ref{ch:do-you-need-this} promised and delivered: your deployment,
yours alone. Removing one is the same deployment and it is not yours alone at
all.

\index{Query type@\code{Query} type!one of the seven is retired in name}%
And one of those seven is already retired in name only, in a graph this book
built from nothing over fifteen chapters, with one author, no other teams and
no clients at all.

\subsection{What That Grows Into}

\index{REST-shaped schema!what the phrase means}%
The shape this converges on is what I call a REST-shaped schema: one where the
root fields are a list of endpoints, each added for one caller, because adding
was the only safe move available. It is not a schema anybody designed. It is
the residue of two years of that decision, and the tell is that the root field
names read like screens rather than like a domain.

\index{REST-shaped schema!the second-order cost}%
The cost is not the size. A large schema is not a problem by itself, and
chapter~\ref{ch:cross-seam-paging} added a whole service to this graph on
purpose. The cost is that a field added for one caller carries no statement
about who else may use it, so within a year it has other callers, and now it
cannot be changed either. Each new field starts free and becomes load-bearing
without anybody deciding that it should.

\index{REST-shaped schema!what it looks like at scale}%
Nikhil Patel, an engineer at Whatnot, published what the far end of this looks
like on a graph much larger than the one in this book: over 2,600 unused
fields, of which nearly 200 were root queries and
mutations~\autocite{patel2025bloat}. His account of why nobody had cleaned
them up is the same argument this chapter has been making, arrived at from
production rather than from a composer:

\begin{quote}
Older versions of our mobile apps may still call queries that have long been
removed from the main branch, and admin tooling dashboards often execute
queries entirely outside our primary repositories. That means static analysis
can't tell the whole story.
\end{quote}

\index{REST-shaped schema!static analysis is the wrong tool}%
Nearly two hundred root fields nobody uses, and the reason they survived is
that the repository does not contain the callers. Reading harder does not help.
Section~\ref{sec:ch16-what-the-composer-is-checking} said the same thing about
the composer and this is the same sentence one layer out: the operation is not
in your source tree, so no tool that reads your source tree will find it.

\subsection{What to Do About It Before It Happens}

\index{REST-shaped schema!the record is the fix}%
The fix is the usage record, and this section exists because a team that does
not have one should know what they are choosing. Everything else on the
subject is mitigation, and I would still do all of it.

\index{schema design!prefer an argument to a new root field}%
Prefer an argument on an existing field to a new root field. A caller who
needs the schedule filtered by track is asking for
\code{sessions(track:)}, and chapter~\ref{ch:cross-seam-paging} put a whole
discovery service behind exactly that instinct rather than adding a
\code{sessionsForTrackPage}. The argument is discoverable by the next caller
who wants the same thing. A field named after one screen is not.

\index{deprecation!name the replacement and then look for the traffic}%
Deprecate with a date, not only a reason.
Chapter~\ref{ch:schema-design} argued that the reason string is the whole user
interface of a deprecation and that it should name the replacement. That
still holds and it is not sufficient, because a reason with no date is a
reason with no plan behind it, and \code{sessionById} is the proof. Put the
month you intend to remove it in the string, and then be honest with yourself
that you will not remove it in that month unless somebody is collecting the
traffic.

\index{schema design!a root field is the most expensive kind}%
And treat a root field as the most expensive thing in the schema. It is
reachable by every client, it is the only position whose parent is
\code{data}, which chapter~\ref{ch:null-blast-radius} measured the cost of,
and once four teams share one \code{Query} type it is the field nobody can
argue about without a meeting. Fields on a type you own are cheap. The shared
root is not.

\index{ownership!the type four services write to and nobody owns}%
Which leaves the word I have been using as though it were settled. Every
recommendation in this section rests on knowing which fields are yours, and
the table two subsections up has four services writing to one \code{Query}
type that belongs to none of them. \code{Session} is worse. Three services
declare it: Sessions holds the rows, Ratings hangs three fields off a key it
can be entered through, and Search points at it with a stub. Ratings owns no
session and can put a field on \code{Session} tomorrow without telling
anybody, and Search could do the same by changing one argument. Composition
will merge all of it without asking whose type it is, and the next chapter is
what happens when somebody has to decide.
\index{schema growth!what a missing usage record does to a schema|)}%
